\chapter{Безмассовый голономный спектр и граница дефектного индекса}

Предыдущая глава закрыла использование тетраэдрической голономии как
источника массы, но сохранила её естественную роль граничной фазы
безмассового переноса. Теперь проверяется точное следствие этой роли:
возникают ли без новых коэффициентов одна инвариантная нулевая ветвь и две
ветви с дробными сдвигами импульса.

\section{Трёхцикл как граничное условие}

На стандартном семейном триплете трёхцикл $C_3$ удовлетворяет
\begin{equation}
 C_3^3=I,
\end{equation}
и имеет собственные значения
\begin{equation}
 1,qquad \omega=e^{2\pi i/3},qquad \bar\omega=e^{-2\pi i/3}.
 \label{eq:v5-c3-eigenvalues}
\end{equation}
Инвариантная собственная линия одномерна и порождена нормированным
вектором
\begin{equation}
 v_0=\frac1{\sqrt3}(1,1,1)^T.
\end{equation}

Рассмотрим один безмассовый хиральный канал на окружности длины $L$ с
оператором
\begin{equation}
 D=-i\frac{d}{ds}
\end{equation}
и граничным условием
\begin{equation}
 \psi(s+L)=C_3\psi(s).
 \label{eq:v5-c3-boundary-condition}
\end{equation}
Спектры операторов Дирака с плоским скручиванием определяются фазами
голономии; это стандартная граница теории плоско скрученных операторов,
см. \path{arXiv:math/0312004}.

\section{Точные дробные ветви}

Для собственного значения $e^{2\pi i\alpha}$ условие
\eqref{eq:v5-c3-boundary-condition} даёт
\begin{equation}
 p_{n,\alpha}=\frac{2\pi}{L}(n+\alpha),
 \qquad n\in\mathbb Z.
\end{equation}
Из \eqref{eq:v5-c3-eigenvalues} следуют три ветви
\begin{align}
 p_{n,-}&=\frac{2\pi}{L}\left(n-\frac13\right),\\
 p_{n,0}&=\frac{2\pi n}{L},\\
 p_{n,+}&=\frac{2\pi}{L}\left(n+\frac13\right).
 \label{eq:v5-third-shift-spectrum}
\end{align}
Следовательно, только инвариантная ветвь имеет нулевой уровень:
\begin{equation}
 p_{0,0}=0,
 \qquad
 \dim_{\mathbb C}\ker D=1.
 \label{eq:v5-one-boundary-zero}
\end{equation}
Минимальный ненулевой модуль импульса равен
\begin{equation}
 \Delta p=\frac{2\pi}{3L}.
 \label{eq:v5-holonomy-third-gap}
\end{equation}
Он является не массой, а голономным сдвигом импульса двух неинвариантных
ветвей.

Машинный аудит перечисляет моды при $|n|\le8$ и находит ровно один нуль.
Инверсия winding заменяет $C_3$ на $C_3^{-1}$ и лишь переставляет ветви
$+1/3$ и $-1/3$.

\begin{theorem}[Параметрически чистый голономный селектор]
На одном хиральном триплетном канале тетраэдрическая голономия $C_3$
порождает без дополнительных коэффициентов спектральные сдвиги
$0,\pm1/3$. Инвариантная собственная линия содержит ровно один комплексный
нулевой уровень, а две сопряжённые ветви отделены щелью
$2\pi/(3L)$.
\end{theorem}

\section{Почему существенна хиральная ветвь
\texorpdfstring{$H_{15}$}{H15}}

Минимальное двухнаправленное блуждание предыдущей главы содержит левый и
правый сдвиги. Если оба считать независимыми, инвариантная семейная линия
даёт два комплексных нулевых состояния --- по одному на направление.

Пакет $H_{15}$ содержит левый нейтринный канал, но не содержит независимого
правого нейтрино. Поэтому ограничение на один хиральный канал имеет
предварительно зафиксированное представительное происхождение и уменьшает
комплексную кратность нуля до единицы.

После дополнительного вещественного условия Майораны постоянный вектор
$v_0$ можно выбрать вещественным. Тогда остаётся одна вещественная нулевая
мода и нечётная чётность
\begin{equation}
 \dim_{\mathbb R}\ker D\equiv1\pmod2.
 \label{eq:v5-boundary-mod-two}
\end{equation}
Но это утверждение условно: оно использует одновременно хиральное
ограничение $H_{15}$ и майорановскую реальность.

\section{Обычный индекс остаётся нулевым}

Наличие одного нулевого уровня нельзя автоматически называть ненулевым
индексом. Оператор на замкнутой окружности действует между одинаковыми
пространствами, а для инвариантной ветви
\begin{equation}
 \dim\ker D=1,
 \qquad
 \dim\operatorname{coker}D=1.
\end{equation}
Поэтому
\begin{equation}
 \operatorname{ind}D=0.
 \label{eq:v5-boundary-index-zero}
\end{equation}

Уравнение \eqref{eq:v5-boundary-mod-two} фиксирует лишь условную
вещественную чётность ядра. Оно ещё не является индексной теоремой для
объёмного вихря или дефектного ядра.

Спектральная асимметрия также не выбирает ориентацию. Для двух дробных
ветвей при $s=0$ получаются вклады $+1/3$ и $-1/3$, которые сокращаются:
\begin{equation}
 \eta_{+}(0)+\eta_{-}(0)=0.
\end{equation}
Знак winding только меняет их местами.

\section{Проверка локализации}

Нулевая собственная функция имеет вид
\begin{equation}
 \psi_0(s)=\frac1{\sqrt L}v_0.
\end{equation}
Она постоянна вдоль всей окружности. На дискретизации из $N$ узлов её
обратный коэффициент участия равен
\begin{equation}
 \sum_{x=1}^{N}|\psi_0(x)|^4=\frac1N.
\end{equation}
Для $N=60$ аудит даёт $1/60$ с машинной точностью. При увеличении системы
эта величина стремится к нулю, что является признаком делокализованной, а
не дефектно-ядерной моды.

\begin{nogocriterion}[Запрет отождествления граничного нуля с частицей]
Один нулевой уровень плоско скрученного оператора нельзя объявлять
локализованным нейтрино или майорановским ядром, пока не построен объёмный
профиль дефекта и не доказана нормируемая локализация поперёк его ядра.
\end{nogocriterion}

\section{Статус ароматных осцилляций}

В собственном базисе $C_3$ три ветви распространяются диагонально. Чтобы
получить наблюдаемую вероятность перехода между ароматами, необходима
независимо выведенная карта между физическим ароматным базисом и
голономными собственными каналами. Текущий проект такую карту не вывел.

Поэтому фиксированные фазы $0,\pm1/3$ являются содержательной спектральной
структурой, но ещё не матрицей нейтринных осцилляций.

\section{Следующий мост к дефектному ядру}

Результат настоящей главы совпадает по кратности с прежним условным
майорановским селектором: одна инвариантная линия и нечётная вещественная
чётность. Различается локализация. Здесь мода граничная и равномерная;
ранняя вихревая ветвь обещала локализованную моду, но принимала
спаривающий конденсат и его профиль как вход.

Следующий гейт должен соединить эти результаты без ручной стенки массы:
проверить, определяют ли уже имеющиеся проекторная суперкривизна $Q(H)$,
кубическая корневая связь и оператор переноса объёмный радиальный профиль,
нормируемую нулевую моду и её ширину без нового коэффициента жёсткости,
масштаба конденсата или выбранной функции $m(r)$.

Если профиль требует хотя бы одного такого входа, граница между
голономным нулём и физическим дефектом остаётся незамкнутой.

\begin{protocol}
\begin{enumerate}
 \item собственные фазы $C_3$:
 \textbf{$0,\pm1/3$};
 \item инвариантная собственная линия:
 \textbf{одномерна};
 \item нулевой уровень одного хирального канала:
 \textbf{ровно один};
 \item минимальный ненулевой сдвиг:
 \textbf{$2\pi/(3L)$};
 \item вещественная чётность Майораны:
 \textbf{нечётная условно на $H_{15}$};
 \item целочисленный индекс:
 \textbf{равен нулю};
 \item суммарная эта-асимметрия:
 \textbf{равна нулю};
 \item локализация нулевой моды:
 \textbf{отсутствует};
 \item физическая матрица осцилляций:
 \textbf{не выведена};
 \item граничный спектральный селектор:
 \textbf{получен};
 \item локализованный нейтринный дефект:
 \textbf{не получен};
 \item физическое замыкание:
 \textbf{не достигнуто}.
\end{enumerate}
\end{protocol}

Машинный сертификат сохранён в
\path{s2t_v5_massless_holonomy_defect_index_gate.py} и
\path{s2t_v5_massless_holonomy_defect_index_gate_results.json}.